\documentclass[10pt]{beamer}

\usetheme{default}
\usecolortheme{default}

\usepackage[utf8]{inputenc}
\usepackage[T2A]{fontenc}
\usepackage[russian]{babel}
\usepackage{amsmath,amssymb,mathtools}
\usepackage{graphicx}
\usepackage{xcolor}
\usepackage{tikz}
\usepackage{booktabs}
\usepackage{array}
\usepackage{hyperref}

% ------------------------------------------------------------
% Безопасная вставка будущих визуальных материалов.
% Если файл пока отсутствует, LaTeX просто пропускает его без ошибки.
% Как только файл появляется по указанному пути, он вставляется автоматически.
% ------------------------------------------------------------

\newcommand{\safeimg}[2][]{%
    \IfFileExists{#2}{\includegraphics[#1]{#2}}{}%
}

\newcommand{\safeanimation}[3][]{%
    \IfFileExists{#2}{%
        \IfFileExists{#3}{%
            \href{run:#3}{\includegraphics[#1]{#2}}%
        }{%
            \includegraphics[#1]{#2}%
        }%
    }{}%
}


% ------------------------------------------------------------
% Цвета и базовый стиль — по MatanPresentation
% ------------------------------------------------------------

\definecolor{defblue}{RGB}{0,0,150}
\definecolor{cardblue}{RGB}{216,216,240}

\newcommand{\blue}[1]{\textcolor{blue}{#1}}
\newcommand{\dblue}[1]{\textcolor{defblue}{#1}}

\setbeamertemplate{navigation symbols}{}
\setbeamertemplate{itemize items}[circle]
\setbeamertemplate{enumerate items}[default]

\setbeamercolor{frametitle}{fg=blue}
\setbeamerfont{frametitle}{size=\large}
\setbeamerfont{title}{size=\Large}
\setbeamerfont{author}{size=\normalsize}

\setbeamersize{
    sidebar width right=1.45cm,
    text margin left=0.48cm,
    text margin right=0.28cm
}

\setbeamertemplate{sidebar right}{
    \begin{tikzpicture}[remember picture,overlay]
        \fill[cardblue]
            (current page.north east)
            rectangle ++(-1.65cm,-\paperheight);

        \node[
            blue,
            anchor=north,
            font=\fontsize{4.8}{5.2}\selectfont\bfseries,
            align=center,
            text width=1.48cm
        ]
        at ([xshift=-0.82cm,yshift=-0.18cm]current page.north east)
        {МатАн. Проект.\\
         Машинное\\
         обучение};

        \node[
            black,
            anchor=north,
            font=\fontsize{4.8}{5.4}\selectfont,
            align=center,
            text width=1.48cm
        ]
        at ([xshift=-0.82cm,yshift=-1.55cm]current page.north east)
        {Вейгандт В.\\
         Зайцев Н.\\
         Горобец О.\\
         Козлов В.\\
         Шагеев Э.};
    \end{tikzpicture}
}

% ------------------------------------------------------------
% Служебные обозначения
% ------------------------------------------------------------

\newcommand{\chapterframe}[2]{%
    \begin{frame}
        \vfill
        \begin{center}
            {\Large \dblue{Глава #1}}\\[0.65em]
            {\LARGE #2}
        \end{center}
        \vfill
    \end{frame}
}

\newcommand{\mlDefinition}{\dblue{Определение}}
\newcommand{\mlStatement}{\dblue{Утверждение}}
\newcommand{\mlRemark}{\dblue{Замечание}}
\newcommand{\mlConsequence}{\dblue{Следствие}}
\newcommand{\mlIdea}{\dblue{Идея}}

% ------------------------------------------------------------
% Метаданные
% ------------------------------------------------------------

\title{Математический анализ в методах машинного обучения}
\author{Вейгандт В., Зайцев Н., Горобец О., Козлов В., Шагеев Э.}
\date{2026}

\begin{document}

% ============================================================
% Титульный слайд
% ============================================================

\begin{frame}
    \titlepage
\end{frame}

% ============================================================
% Глава 1
% ============================================================

\chapterframe{1}{Постановка задачи машинного обучения}

\begin{frame}{Обучение с учителем}
    \mlDefinition

    Пусть $\mathcal X$ --- множество объектов, $\mathcal Y$ --- множество допустимых ответов.
    \textit{Обучающей выборкой} называется конечный набор
    \[
        D=\{(x_i,y_i)\}_{i=1}^{n},
        \qquad x_i\in\mathcal X,\quad y_i\in\mathcal Y.
    \]

    \mlDefinition

    В задаче \textit{supervised learning} (обучения с учителем) по выборке $D$
    требуется построить алгоритм
    \[
        f:\mathcal X\to\mathcal Y,
    \]
    который по описанию нового объекта $x$ предсказывает соответствующий ответ.

    \mlRemark

    Правильные ответы $y_i$ известны только для обучающих объектов; качество модели
    определяется её способностью переносить найденную закономерность на новые данные.

    \begin{center}
        \safeimg[width=0.72\linewidth]{img/supervised_learning_scheme.pdf}
    \end{center}
\end{frame}

\begin{frame}{Признаковое описание задачи}
    \mlDefinition

    Числовым \textit{признаковым описанием} объекта называется вектор
    \[
        x=(x_1,\ldots,x_d)^T\in\mathbb R^d,
    \]
    где координаты $x_1,\ldots,x_d$ называются признаками объекта.

    Таким образом, обучающая выборка имеет вид
    \[
        D=\{(x_i,y_i)\}_{i=1}^{n},
        \qquad x_i\in\mathbb R^d.
    \]

    \mlDefinition

    \begin{itemize}
        \item если $\mathcal Y=\mathbb R$, рассматривается задача \textit{регрессии};
        \item если $\mathcal Y=\{1,\ldots,K\}$, рассматривается задача \textit{классификации}.
    \end{itemize}

    \mlRemark

    После перехода к признакам каждый объект соответствует точке в $d$-мерном
    пространстве, а обучение сводится к поиску зависимости между координатами этой точки
    и ответом $y$.
\end{frame}

\begin{frame}{Модель как параметризованное семейство}
    \mlDefinition

    \textit{Параметрической моделью} называется семейство отображений
    \[
        \mathcal F=
        \{f_\theta:\mathcal X\to\mathcal Y\mid \theta\in\Theta\},
    \]
    где $\Theta$ --- пространство параметров, а $\theta$ --- вектор параметров модели.

    Если $\Theta\subseteq\mathbb R^m$, то
    \[
        \theta=(\theta_1,\ldots,\theta_m)^T.
    \]

    \textbf{Пример.}
    Для семейства прямых
    \[
        f_{w,b}(x)=wx+b
    \]
    пространство параметров равно $\Theta=\mathbb R^2$, а $\theta=(w,b)^T$.

    \mlRemark

    До обучения задано всё семейство $\mathcal F$; результат обучения --- конкретная
    функция $f_{\hat\theta}$ с найденным набором параметров $\hat\theta$.
\end{frame}

\begin{frame}{Функция потерь и эмпирический риск}
    \mlDefinition

    \textit{Функцией потерь} называется отображение
    \[
        \ell:\mathcal Y\times\mathcal Y\to\mathbb R_{\ge 0},
    \]
    где $\ell(y,\hat y)$ измеряет ошибку одного предсказания $\hat y$ при истинном ответе $y$.

    \mlDefinition

    \textit{Эмпирическим риском} модели $f_\theta$ на выборке $D$ называется
    \[
        L_D(\theta)
        =\frac1n\sum_{i=1}^{n}
        \ell\bigl(y_i,f_\theta(x_i)\bigr).
    \]

    После фиксации выборки $D$ величина $L_D$ зависит только от параметров модели.
    Если $\Theta\subseteq\mathbb R^m$, то
    \[
        \boxed{L_D:\mathbb R^m\to\mathbb R}
    \]
    является функцией нескольких переменных.

    \mlRemark

    Именно эта функция связывает постановку задачи машинного обучения с задачей
    исследования и минимизации функций нескольких переменных.
\end{frame}

\begin{frame}{Обучение модели}
    \mlDefinition

    Обучением параметрической модели называется задача поиска параметров
    \[
        \boxed{
        \hat\theta\in
        \operatorname*{arg\,min}_{\theta\in\Theta}L_D(\theta)
        }.
    \]

    В развёрнутом виде:
    \[
        \hat\theta\in
        \operatorname*{arg\,min}_{\theta\in\Theta}
        \frac1n\sum_{i=1}^{n}
        \ell\bigl(y_i,f_\theta(x_i)\bigr).
    \]

    \mlConsequence

    При $\Theta\subseteq\mathbb R^m$ задача обучения становится задачей оптимизации
    функции нескольких переменных $L_D(\theta)$.

    \mlRemark

    Разные алгоритмы машинного обучения отличаются выбором семейства $\mathcal F$,
    функции потерь $\ell$ и метода решения полученной задачи оптимизации.

    % Manim: animations/01_training_process.mp4
    \begin{center}
        \href{run:animations/01_training_process.mp4}{%
        \safeimg[width=0.68\linewidth]{img/01_training_process_poster.png}}
    \end{center}
\end{frame}

\begin{frame}{Обучающая, валидационная и тестовая выборки}
    Для оценки способности модели работать на новых данных исходный набор объектов
    разделяют на непересекающиеся части:
    \[
        D=D_{\mathrm{train}}\sqcup
        D_{\mathrm{validation}}\sqcup
        D_{\mathrm{test}}.
    \]

    \begin{itemize}
        \item $D_{\mathrm{train}}$ используется для определения параметров $\hat\theta$;
        \item $D_{\mathrm{validation}}$ --- для выбора гиперпараметров и сравнения вариантов модели;
        \item $D_{\mathrm{test}}$ --- для итоговой оценки качества после завершения обучения.
    \end{itemize}

    \mlDefinition

    Ошибка на обучающей выборке:
    \[
        L_{\mathrm{train}}(\hat\theta).
    \]
    Ошибка на тестовой выборке определяется той же функцией потерь, но на объектах,
    не участвовавших в подборе параметров.

    \mlRemark

    Малое значение $L_{\mathrm{train}}$ не гарантирует малого значения тестовой ошибки:
    модель должна не только описывать обучающие объекты, но и \textit{обобщать} найденную зависимость.

    \begin{center}
        \safeimg[width=0.76\linewidth]{img/train_val_test.pdf}
    \end{center}
\end{frame}

% ============================================================
% Глава 2
% ============================================================

\chapterframe{2}{Линейные модели}

\begin{frame}{Linear Regression (линейная регрессия)}
    \mlDefinition

    \textit{Линейной регрессионной моделью} называется функция
    \[
        f_\theta(x)=x^T\theta,
        \qquad x,\theta\in\mathbb R^d.
    \]

    Свободный коэффициент можно включить в вектор параметров,
    добавив к признаковому вектору единичную координату.

    Для одного признака получаем
    \[
        \hat y=f_{w,b}(x)=wx+b.
    \]

    \mlRemark

    Слово «линейная» относится к зависимости от параметров.
    Например,
    \[
        \hat y=\theta_0+\theta_1x+\theta_2x^2
    \]
    остаётся линейной моделью по параметрам $\theta_0,\theta_1,\theta_2$.

    \begin{center}
        \safeimg[width=0.72\linewidth]{img/linear_regression.pdf}
    \end{center}
\end{frame}

\begin{frame}{Mean Squared Error (среднеквадратичная ошибка)}
    Для объекта $(x_i,y_i)$ определим остаток
    \[
        r_i=y_i-f_\theta(x_i).
    \]

    \mlDefinition

    Квадратичная функция потерь:
    \[
        \ell(y,\hat y)=(y-\hat y)^2.
    \]
    Соответствующий эмпирический риск имеет вид
    \[
        L(\theta)=
        \frac1n\sum_{i=1}^{n}
        \bigl(y_i-x_i^T\theta\bigr)^2.
    \]

    Если строки матрицы $X\in\mathbb R^{n\times d}$ равны $x_i^T$, а
    $Y=(y_1,\ldots,y_n)^T$, то
    \[
        \boxed{L(\theta)=\frac1n\|Y-X\theta\|_2^2}.
    \]

    \mlRemark

    Возведение остатка в квадрат исключает взаимную компенсацию положительных и
    отрицательных ошибок и сильнее штрафует большие отклонения.
\end{frame}

\begin{frame}{MSE как функция нескольких переменных}
    Для модели одного признака
    \[
        f_{w,b}(x)=wx+b
    \]
    функция потерь принимает вид
    \[
        L(w,b)=
        \frac1n\sum_{i=1}^{n}(y_i-wx_i-b)^2.
    \]

    Следовательно,
    \[
        L:\mathbb R^2\to\mathbb R.
    \]

    Множества
    \[
        \{(w,b):L(w,b)=C\}
    \]
    являются линиями уровня функции потерь.

    \begin{columns}[T,onlytextwidth]
        \column{0.50\textwidth}
        \centering
        \safeimg[width=0.96\linewidth]{img/mse_surface.pdf}

        \column{0.50\textwidth}
        \centering
        \safeimg[width=0.96\linewidth]{img/mse_contours.pdf}
    \end{columns}

    \mlConsequence\;
    Минимум поверхности соответствует параметрам наилучшей прямой в смысле MSE.
\end{frame}

\begin{frame}{Градиент MSE}
    \small
    \mlStatement

    Для
    \[
        L(\theta)=\frac1n\|Y-X\theta\|_2^2
    \]
    градиент по параметрам равен
    \[
        \boxed{
        \nabla_\theta L(\theta)
        =\frac{2}{n}X^T(X\theta-Y)
        }.
    \]

    В координатной форме
    \[
        \nabla L(\theta)=
        \begin{pmatrix}
            \dfrac{\partial L}{\partial\theta_1}\\
            \vdots\\
            \dfrac{\partial L}{\partial\theta_d}
        \end{pmatrix}.
    \]

    \mlRemark

    Частная производная $\partial L/\partial\theta_j$ характеризует локальную
    чувствительность ошибки к изменению $j$-го параметра при фиксированных остальных.
    Градиент объединяет эти величины в один вектор.

    % Manim: animations/02_linear_regression_loss.mp4
    \begin{center}
        \href{run:animations/02_linear_regression_loss.mp4}{%
        \safeimg[width=0.64\linewidth]{img/02_linear_regression_loss_poster.png}}
    \end{center}
\end{frame}

\begin{frame}{Logistic Regression (логистическая регрессия)}
    Рассмотрим бинарную классификацию:
    \[
        y\in\{0,1\}.
    \]

    Линейный score
    \[
        z=x^T\theta
    \]
    принимает произвольные вещественные значения и потому не может непосредственно
    интерпретироваться как вероятность.

    \mlDefinition

    В \textit{Logistic Regression} вероятность класса $1$ задаётся
    \[
        p_\theta(x)=P_\theta(Y=1\mid x)=\sigma(x^T\theta),
    \]
    где
    \[
        \boxed{\sigma(z)=\frac{1}{1+e^{-z}}}.
    \]

    При стандартном пороге $1/2$
    \[
        \hat y=\mathbb I\{p_\theta(x)\ge 1/2\}.
    \]

    \begin{columns}[T,onlytextwidth]
        \column{0.50\textwidth}
        \centering
        \safeimg[width=0.95\linewidth]{img/sigmoid.pdf}
        \column{0.50\textwidth}
        \centering
        \safeimg[width=0.95\linewidth]{img/logistic_boundary.pdf}
    \end{columns}
\end{frame}

\begin{frame}{Maximum Likelihood и Cross-Entropy}
    Пусть
    \[
        p_i=p_\theta(x_i)=P_\theta(Y_i=1\mid x_i).
    \]
    Тогда вероятность наблюдать ответ $y_i\in\{0,1\}$ равна
    \[
        p_i^{y_i}(1-p_i)^{1-y_i}.
    \]

    При условной независимости наблюдений функция правдоподобия:
    \[
        \mathcal L(\theta)
        =\prod_{i=1}^{n}p_i^{y_i}(1-p_i)^{1-y_i}.
    \]

    \mlStatement

    Максимизация $\mathcal L(\theta)$ эквивалентна минимизации отрицательного
    среднего лог-правдоподобия:
    \[
        \boxed{
        L(\theta)=
        -\frac1n\sum_{i=1}^{n}
        \left[y_i\log p_i+(1-y_i)\log(1-p_i)\right]
        }.
    \]

    \mlDefinition\;
    Эта функция называется \textit{binary cross-entropy} (бинарной кросс-энтропией).

    % Manim: animations/03_logistic_sigmoid.mp4
    \begin{center}
        \href{run:animations/03_logistic_sigmoid.mp4}{%
        \safeimg[width=0.56\linewidth]{img/03_logistic_sigmoid_poster.png}}
    \end{center}
\end{frame}

% ============================================================
% Глава 3
% ============================================================

\chapterframe{3}{Градиентные методы оптимизации}

\begin{frame}{Геометрический смысл градиента}
    Пусть
    \[
        L:\mathbb R^d\to\mathbb R
    \]
    дифференцируема в точке $\theta$.

    \mlStatement

    Производная $L$ по единичному направлению $v$, $\|v\|=1$, равна
    \[
        D_vL(\theta)=\langle\nabla L(\theta),v\rangle.
    \]

    Из неравенства Коши--Буняковского
    \[
        -\|\nabla L(\theta)\|
        \le D_vL(\theta)
        \le \|\nabla L(\theta)\|.
    \]

    Наибольшее локальное убывание достигается в направлении
    \[
        \boxed{
        v=-\frac{\nabla L(\theta)}{\|\nabla L(\theta)\|}
        }.
    \]

    \mlRemark\;
    Градиент ортогонален линии уровня в регулярной точке функции.

    \begin{center}
        \safeimg[width=0.65\linewidth]{img/gradient_geometry.pdf}
    \end{center}
\end{frame}

\begin{frame}{Gradient Descent (градиентный спуск)}
    \mlDefinition

    \textit{Gradient Descent} задаётся последовательностью
    \[
        \boxed{
        \theta_{t+1}
        =\theta_t-\eta_t\nabla L(\theta_t)
        },
    \]
    где $\eta_t>0$ --- \textit{learning rate} (скорость обучения).

    На каждой итерации выполняются два шага:
    \begin{enumerate}
        \item вычисляется градиент $\nabla L(\theta_t)$;
        \item параметры изменяются в направлении антиградиента.
    \end{enumerate}

    \mlRemark

    Метод использует только локальную информацию о функции $L$ и не требует явного
    аналитического выражения точки минимума.

    \mlConsequence

    При удачном выборе $\eta_t$ последовательность $L(\theta_t)$ должна уменьшаться,
    а параметры $\theta_t$ --- приближаться к стационарной точке.

    % Manim: animations/04_gradient_descent.mp4
    \begin{center}
        \href{run:animations/04_gradient_descent.mp4}{%
        \safeimg[width=0.64\linewidth]{img/04_gradient_descent_poster.png}}
    \end{center}
\end{frame}

\begin{frame}{Влияние learning rate}
    Значение $\eta$ определяет длину шага вдоль антиградиента.

    Для одномерной функции
    \[
        f(u)=u^2
    \]
    итерация Gradient Descent принимает вид
    \[
        u_{t+1}=u_t-2\eta u_t=(1-2\eta)u_t.
    \]

    Отсюда непосредственно видно, что характер последовательности зависит от $\eta$.

    \begin{itemize}
        \item малое $\eta$ --- устойчивое, но медленное уменьшение ошибки;
        \item умеренное $\eta$ --- быстрая сходимость;
        \item большое $\eta$ --- колебания около минимума;
        \item чрезмерно большое $\eta$ --- расходимость.
    \end{itemize}

    \begin{columns}[T,onlytextwidth]
        \column{0.50\textwidth}
        \centering
        \safeimg[width=0.95\linewidth]{img/learning_rate.pdf}
        \column{0.50\textwidth}
        \centering
        % Manim: animations/05_learning_rate.mp4
        \safeanimation[width=0.95\linewidth]{img/05_learning_rate_poster.png}{animations/05_learning_rate.mp4}
    \end{columns}
\end{frame}

\begin{frame}{Momentum и Nesterov}
    Пусть $g_t=\nabla L(\theta_t)$.

    \mlDefinition\; \textbf{Momentum (метод момента)}
    \[
        v_t=\beta v_{t-1}+g_t,
        \qquad
        \theta_{t+1}=\theta_t-\eta v_t,
        \qquad 0\le\beta<1.
    \]

    Вектор $v_t$ аккумулирует информацию о предыдущих направлениях движения.
    Это уменьшает поперечные колебания в вытянутых долинах функции потерь.

    \mlDefinition\; \textbf{Nesterov Momentum (метод Нестерова)}
    \[
        g_t=
        \nabla L\!\left(\theta_t-\eta\beta v_{t-1}\right),
    \]
    \[
        v_t=\beta v_{t-1}+g_t,
        \qquad
        \theta_{t+1}=\theta_t-\eta v_t.
    \]

    \mlRemark\;
    В методе Нестерова градиент вычисляется в точке, учитывающей предполагаемое
    движение по накопленному импульсу.

    \begin{center}
        \safeimg[width=0.62\linewidth]{img/momentum_valley.pdf}
    \end{center}
\end{frame}

\begin{frame}{AdaGrad и RMSProp}
    Пусть $g_t=\nabla L(\theta_t)$, а $\odot$ обозначает покоординатное умножение.

    \mlDefinition\; \textbf{AdaGrad}
    \[
        G_t=G_{t-1}+g_t\odot g_t,
    \]
    \[
        \theta_{t+1}
        =\theta_t-
        \eta\frac{g_t}{\sqrt{G_t}+\varepsilon}.
    \]

    Эффективный размер шага автоматически уменьшается по координатам, где накоплены
    большие значения градиента.

    \mlDefinition\; \textbf{RMSProp}
    \[
        s_t=\rho s_{t-1}+(1-\rho)g_t\odot g_t,
    \]
    \[
        \theta_{t+1}
        =\theta_t-
        \eta\frac{g_t}{\sqrt{s_t}+\varepsilon}.
    \]

    \mlRemark

    В AdaGrad сумма квадратов градиентов монотонно растёт, поэтому шаг может стать
    слишком малым. RMSProp заменяет полное накопление экспоненциальным скользящим средним.
\end{frame}

\begin{frame}{Adam}
    \mlDefinition

    \textit{Adam} использует экспоненциальные средние первого и второго моментов градиента:
    \[
        m_t=\beta_1m_{t-1}+(1-\beta_1)g_t,
    \]
    \[
        v_t=\beta_2v_{t-1}+(1-\beta_2)g_t\odot g_t.
    \]

    Для компенсации начальных значений $m_0=v_0=0$ применяются поправки
    \[
        \hat m_t=\frac{m_t}{1-\beta_1^t},
        \qquad
        \hat v_t=\frac{v_t}{1-\beta_2^t}.
    \]

    Итоговое обновление:
    \[
        \boxed{
        \theta_{t+1}
        =\theta_t-
        \eta\frac{\hat m_t}{\sqrt{\hat v_t}+\varepsilon}
        }.
    \]

    \mlRemark

    Первый момент задаёт сглаженное направление движения, второй --- адаптивный
    масштаб шага по отдельным координатам.
\end{frame}

\begin{frame}{Сравнение градиентных методов}
    Для корректного сравнения оптимизаторы запускаются
    \begin{itemize}
        \item на одной функции потерь;
        \item из одной стартовой точки;
        \item при фиксированных гиперпараметрах каждого метода.
    \end{itemize}

    Сравниваются
    \[
        \text{GD},\quad
        \text{Momentum},\quad
        \text{Nesterov},\quad
        \text{AdaGrad},\quad
        \text{RMSProp},\quad
        \text{Adam}.
    \]

    Основные характеристики:
    \begin{itemize}
        \item траектория параметров на линиях уровня;
        \item скорость уменьшения $L(\theta_t)$;
        \item число итераций до заданной точности;
        \item устойчивость к различному масштабу координат.
    \end{itemize}

    \begin{columns}[T,onlytextwidth]
        \column{0.50\textwidth}
        \centering
        \safeimg[width=0.95\linewidth]{img/optimizer_trajectories.pdf}
        \column{0.50\textwidth}
        \centering
        \safeimg[width=0.95\linewidth]{img/optimizer_loss.pdf}
    \end{columns}

    % Manim: animations/06_optimizers.mp4
\end{frame}

% ============================================================
% Глава 4
% ============================================================

\chapterframe{4}{Решающие деревья и Gradient Boosting}

\begin{frame}{Decision Tree (решающее дерево)}
    \mlDefinition

    Внутренняя вершина решающего дерева содержит предикат
    \[
        x_j\le t,
    \]
    где $j$ --- номер признака, $t$ --- порог.

    Для множества объектов $S$ предикат определяет разбиение
    \[
        S_L(j,t)=\{(x_i,y_i)\in S:x_{ij}\le t\},
    \]
    \[
        S_R(j,t)=\{(x_i,y_i)\in S:x_{ij}>t\}.
    \]

    Внутренние вершины задают последовательность условий, а листья содержат итоговое
    предсказание модели.

    \mlRemark

    В отличие от линейных моделей, структура дерева сама является частью решения:
    необходимо выбрать признаки, пороги и глубину разбиений.

    \begin{columns}[T,onlytextwidth]
        \column{0.50\textwidth}
        \centering
        \safeimg[width=0.94\linewidth]{img/tree_split.pdf}
        \column{0.50\textwidth}
        \centering
        \safeimg[width=0.94\linewidth]{img/decision_tree.pdf}
    \end{columns}
\end{frame}

\begin{frame}{Критерий разбиения}
    Пусть $H(S)$ --- мера неоднородности ответов объектов, попавших в вершину.

    \mlDefinition

    \textit{Качеством разбиения} $(j,t)$ назовём уменьшение неоднородности
    \[
        Q(S,j,t)
        =H(S)
        -\frac{|S_L|}{|S|}H(S_L)
        -\frac{|S_R|}{|S|}H(S_R).
    \]

    На каждом шаге выбирается
    \[
        \boxed{
        (j^\ast,t^\ast)
        \in
        \operatorname*{arg\,max}_{j,t}Q(S,j,t)
        }.
    \]

    \mlRemark

    Это локальная задача оптимизации: алгоритм не ищет сразу оптимальное дерево целиком,
    а выбирает наилучшее разбиение в текущей вершине.

    \mlConsequence

    Такое построение называется \textit{жадным}: каждый шаг оптимален относительно
    текущей вершины, но не обязан давать глобально оптимальную структуру дерева.

    \begin{center}
        \safeimg[width=0.68\linewidth]{img/tree_candidates.pdf}
    \end{center}
\end{frame}

\begin{frame}{Критерии для регрессии и классификации}
    \small
    \dblue{Регрессия}

    В качестве неоднородности можно использовать дисперсию ответов в вершине:
    \[
        H_{\mathrm{MSE}}(S)
        =\frac1{|S|}
        \sum_{(x_i,y_i)\in S}(y_i-\bar y_S)^2,
    \]
    где
    \[
        \bar y_S=\frac1{|S|}\sum_{(x_i,y_i)\in S}y_i.
    \]
    Предсказанием листа обычно является $\bar y_S$.

    \dblue{Классификация}

    Критерий Джини:
    \[
        G(S)=1-\sum_{k=1}^{K}p_k^2.
    \]
    Энтропия:
    \[
        H(S)=-\sum_{k=1}^{K}p_k\log p_k.
    \]

    Здесь $p_k$ --- доля объектов класса $k$ в вершине.

    \mlRemark

    Хороший split уменьшает неоднородность ответов в дочерних вершинах
    по сравнению с исходной вершиной.
\end{frame}

\begin{frame}{Жадное построение дерева}
    \mlStatement

    Полный перебор всех возможных структур дерева практически неосуществим, поэтому
    дерево строится рекурсивно и жадно.

    \begin{enumerate}
        \item В текущую вершину приходит множество объектов $S$.
        \item Перебираются допустимые признаки $j$ и пороги $t$.
        \item Для каждого кандидата вычисляется $Q(S,j,t)$.
        \item Выбирается локально лучшее разбиение $(j^\ast,t^\ast)$.
        \item Процедура повторяется в дочерних вершинах.
    \end{enumerate}

    Рост дерева ограничивают, например,
    \begin{itemize}
        \item максимальной глубиной;
        \item минимальным числом объектов в вершине;
        \item минимальным улучшением критерия.
    \end{itemize}

    \mlRemark

    Параметры дерева включают дискретный выбор признака и структуры, поэтому производная
    по этим величинам в обычном смысле не определена. Gradient Descent непосредственно
    к построению дерева не применяется.

    % Manim: animations/07_decision_tree.mp4
    \begin{center}
        \href{run:animations/07_decision_tree.mp4}{%
        \safeimg[width=0.58\linewidth]{img/07_decision_tree_poster.png}}
    \end{center}
\end{frame}

\begin{frame}{Random Forest и Gradient Boosting}
    \begin{columns}[T,onlytextwidth]
        \column{0.49\textwidth}
        \dblue{Random Forest (случайный лес)}

        Строится множество деревьев
        \[
            h_1,\ldots,h_M,
        \]
        после чего их ответы агрегируются усреднением или голосованием.

        Для уменьшения корреляции между деревьями используются
        \begin{itemize}
            \item bootstrap-выборки объектов;
            \item случайные подмножества признаков при выборе split'а.
        \end{itemize}

        \column{0.49\textwidth}
        \dblue{Gradient Boosting (градиентный бустинг)}

        Строится последовательность композиций
        \[
            F_0\to F_1\to\cdots\to F_M.
        \]
        Каждое следующее дерево зависит от уже построенного ансамбля и исправляет его
        текущую ошибку.
    \end{columns}

    \mlRemark

    Random Forest строит разнообразные деревья преимущественно независимо;
    Gradient Boosting строит деревья последовательно.

    \begin{center}
        \safeimg[width=0.72\linewidth]{img/rf_vs_boosting.pdf}
    \end{center}
\end{frame}

\begin{frame}{Gradient Boosting (градиентный бустинг)}
    Пусть после $m-1$ шагов построена композиция $F_{m-1}$.

    \mlDefinition

    Следующая модель добавляется по правилу
    \[
        \boxed{
        F_m(x)=F_{m-1}(x)+\eta h_m(x)
        },
    \]
    где $h_m$ --- новое дерево, а $\eta>0$ --- шаг бустинга.

    Для дифференцируемой функции потерь вводятся \textit{псевдоостатки}
    \[
        \boxed{
        r_i^{(m)}
        =-
        \left.
        \frac{\partial\ell(y_i,F(x_i))}{\partial F(x_i)}
        \right|_{F=F_{m-1}}
        }.
    \]

    Новое дерево обучается приближать значения $r_i^{(m)}$ на обучающих объектах.

    \mlRemark

    Величина $\eta$ играет роль шага: малое значение делает обновления более осторожными,
    но требует большего числа деревьев.
\end{frame}

\begin{frame}{Gradient Boosting как спуск в пространстве функций}
    Для квадратичной функции потерь
    \[
        \ell(y,F)=(y-F)^2
    \]
    имеем
    \[
        -\frac{\partial\ell(y,F)}{\partial F}
        =2(y-F).
    \]
    Следовательно, отрицательный градиент пропорционален обычному остатку
    \[
        y_i-F_{m-1}(x_i).
    \]

    \mlConsequence

    При MSE очередное дерево обучается приближать ошибки текущего ансамбля.

    \mlRemark

    В Gradient Descent изменяется конечномерный вектор параметров $\theta$;
    в Gradient Boosting последовательно изменяется функция $F$ посредством добавления
    новых базовых моделей $h_m$.

    \begin{columns}[T,onlytextwidth]
        \column{0.50\textwidth}
        \centering
        \safeimg[width=0.95\linewidth]{img/boosting_steps.pdf}
        \column{0.50\textwidth}
        \centering
        \safeimg[width=0.95\linewidth]{img/boosting_loss.pdf}
    \end{columns}

    % Manim: animations/08_gradient_boosting.mp4
\end{frame}

% ============================================================
% Глава 5
% ============================================================

\chapterframe{5}{Дополнительные методы}

\begin{frame}{Principal Component Analysis (метод главных компонент)}
    Пусть строки матрицы $X$ содержат центрированные объекты, а
    \[
        C=\frac1N X^TX
    \]
    --- ковариационная матрица.

    Для единичного направления $w$ проекция объекта $x$ равна $w^Tx$, а дисперсия
    проекций определяется величиной $w^TCw$.

    \mlDefinition

    Первая главная компонента задаётся задачей
    \[
        \boxed{
        \max_{\|w\|=1}w^TCw
        }.
    \]

    \mlStatement

    Решением является собственный вектор матрицы $C$, соответствующий наибольшему
    собственному значению.

    \mlRemark

    PCA выбирает ортогональные направления максимальной дисперсии и позволяет оставить
    первые компоненты, отбросив направления с относительно малой вариативностью.

    \begin{columns}[T,onlytextwidth]
        \column{0.50\textwidth}
        \centering
        \safeimg[width=0.94\linewidth]{img/pca.pdf}
        \column{0.50\textwidth}
        \centering
        % Manim: animations/09_pca.mp4
        \safeanimation[width=0.94\linewidth]{img/09_pca_poster.png}{animations/09_pca.mp4}
    \end{columns}
\end{frame}

\begin{frame}{Monte Carlo (метод Монте-Карло)}
    Рассмотрим интеграл
    \[
        I=\int_a^b f(x)\,dx.
    \]
    Если $X\sim U(a,b)$, то
    \[
        \mathbb E[f(X)]
        =\frac{1}{b-a}\int_a^b f(x)\,dx,
    \]
    поэтому
    \[
        I=(b-a)\mathbb E[f(X)].
    \]

    \mlDefinition

    Для независимых $X_1,\ldots,X_N\sim U(a,b)$ оценка Монте-Карло имеет вид
    \[
        \boxed{
        \hat I_N=
        \frac{b-a}{N}\sum_{i=1}^{N}f(X_i)
        }.
    \]

    \mlRemark

    В основе метода лежит замена математического ожидания выборочным средним.
    При увеличении $N$ оценка стабилизируется в соответствии с законом больших чисел.

    \begin{columns}[T,onlytextwidth]
        \column{0.50\textwidth}
        \centering
        \safeimg[width=0.94\linewidth]{img/monte_carlo.pdf}
        \column{0.50\textwidth}
        \centering
        % Manim: animations/10_monte_carlo.mp4
        \safeanimation[width=0.94\linewidth]{img/10_monte_carlo_poster.png}{animations/10_monte_carlo.mp4}
    \end{columns}
\end{frame}

\begin{frame}{Ant Colony Optimization (муравьиный алгоритм)}
    Рассмотрим граф, в котором из вершины $i$ допустимы переходы в вершины $j\in A_i$.

    \mlDefinition

    Вероятность перехода по ребру $(i,j)$ задаётся
    \[
        P_{ij}
        =
        \frac{\tau_{ij}^{\alpha}\eta_{ij}^{\beta}}
        {\sum\limits_{k\in A_i}\tau_{ik}^{\alpha}\eta_{ik}^{\beta}},
    \]
    где $\tau_{ij}$ --- накопленный феромон, $\eta_{ij}$ --- локальная привлекательность
    перехода, а $\alpha,\beta$ задают их относительное влияние.

    После итерации феромон обновляется:
    \[
        \boxed{
        \tau_{ij}\leftarrow
        (1-\rho)\tau_{ij}+\Delta\tau_{ij}
        }.
    \]

    \mlRemark

    Испарение $(1-\rho)\tau_{ij}$ препятствует бесконечному накоплению старой информации,
    а $\Delta\tau_{ij}$ усиливает рёбра, участвовавшие в более качественных решениях.

    \begin{columns}[T,onlytextwidth]
        \column{0.33\textwidth}
        \centering
        \safeimg[width=0.94\linewidth]{img/ant_colony_early.pdf}
        \column{0.33\textwidth}
        \centering
        \safeimg[width=0.94\linewidth]{img/ant_colony_late.pdf}
        \column{0.33\textwidth}
        \centering
        % Manim: animations/11_ant_colony.mp4
        \safeanimation[width=0.94\linewidth]{img/11_ant_colony_poster.png}{animations/11_ant_colony.mp4}
    \end{columns}
\end{frame}

% ============================================================
% Глава 6
% ============================================================

\chapterframe{6}{Практическая часть}

\begin{frame}{Программная реализация}
    Практическая часть реализуется на языке Python.

    \begin{itemize}
        \item \texttt{NumPy} --- матричные и векторные вычисления;
        \item \texttt{Matplotlib} --- построение статических графиков;
        \item \texttt{Manim} --- математические анимации;
        \item \texttt{scikit-learn} --- дополнительная проверка полученных результатов.
    \end{itemize}

    \mlStatement

    Основные алгоритмы реализуются самостоятельно, чтобы контролировать
    \begin{itemize}
        \item вычисление градиентов и обновление параметров;
        \item выбор split'ов решающего дерева;
        \item построение последовательности деревьев в Gradient Boosting.
    \end{itemize}

    Для сравнения методов фиксируются исходные данные, стартовые точки и параметры
    эксперимента.

    \mlRemark

    Такой подход позволяет сопоставлять не только итоговые значения loss, но и сами
    траектории и промежуточные состояния алгоритмов.

    \vspace{0.25cm}
    \centering
    \texttt{PLACEHOLDER: ссылка на репозиторий}
\end{frame}

\begin{frame}{Эксперимент: Linear Regression}
    Рассматривается синтетическая выборка одного признака и модель
    \[
        \hat y=wx+b.
    \]

    Функция потерь:
    \[
        L(w,b)=
        \frac1n\sum_{i=1}^{n}(y_i-wx_i-b)^2.
    \]

    В эксперименте фиксируются
    \begin{itemize}
        \item поверхность и линии уровня $L(w,b)$;
        \item начальная точка $(w_0,b_0)$;
        \item последовательность параметров $(w_t,b_t)$;
        \item изменение прямой $wx+b$ в процессе оптимизации.
    \end{itemize}

    \begin{columns}[T,onlytextwidth]
        \column{0.50\textwidth}
        \centering
        \safeimg[width=0.95\linewidth]{img/linear_regression.pdf}
        \column{0.50\textwidth}
        \centering
        \safeimg[width=0.95\linewidth]{img/mse_contours.pdf}
    \end{columns}

    % Manim: animations/02_linear_regression_loss.mp4
\end{frame}

\begin{frame}{Эксперимент: сравнение оптимизаторов}
    На одной двумерной функции из одной начальной точки запускаются
    \[
        \text{GD},\quad
        \text{Momentum},\quad
        \text{Nesterov},\quad
        \text{AdaGrad},\quad
        \text{RMSProp},\quad
        \text{Adam}.
    \]

    Для каждого метода сохраняются
    \begin{itemize}
        \item последовательность параметров $\theta_t$;
        \item значения $L(\theta_t)$;
        \item число итераций;
        \item расстояние до известной точки минимума.
    \end{itemize}

    \mlRemark

    Сравнение траекторий показывает геометрическое различие методов, а график
    $L(\theta_t)$ позволяет сравнить скорость уменьшения функции потерь.

    \begin{columns}[T,onlytextwidth]
        \column{0.50\textwidth}
        \centering
        \safeimg[width=0.95\linewidth]{img/optimizer_trajectories.pdf}
        \column{0.50\textwidth}
        \centering
        \safeimg[width=0.95\linewidth]{img/optimizer_loss.pdf}
    \end{columns}

    % Manim: animations/06_optimizers.mp4
\end{frame}

\begin{frame}{Эксперимент: Decision Tree}
    На небольшом двумерном наборе данных вручную воспроизводятся первые шаги
    жадного построения дерева.

    Для текущей вершины $S$ вычисляются значения
    \[
        Q(S,j,t)
    \]
    для допустимых признаков и порогов, после чего выбирается
    \[
        (j^\ast,t^\ast)
        \in\operatorname*{arg\,max}_{j,t}Q(S,j,t).
    \]

    В эксперименте сохраняются
    \begin{itemize}
        \item кандидатные split'ы;
        \item значения критерия для каждого кандидата;
        \item выбранная граница;
        \item итоговая структура дерева.
    \end{itemize}

    \begin{columns}[T,onlytextwidth]
        \column{0.50\textwidth}
        \centering
        \safeimg[width=0.95\linewidth]{img/tree_candidates.pdf}
        \column{0.50\textwidth}
        \centering
        \safeimg[width=0.95\linewidth]{img/decision_tree.pdf}
    \end{columns}

    % Manim: animations/07_decision_tree.mp4
\end{frame}

\begin{frame}{Эксперимент: Gradient Boosting}
    Строится последовательность композиций
    \[
        F_0,F_1,\ldots,F_M,
        \qquad
        F_m=F_{m-1}+\eta h_m.
    \]

    На каждом шаге сохраняются
    \begin{itemize}
        \item текущие предсказания $F_{m-1}(x_i)$;
        \item residuals или pseudo-residuals $r_i^{(m)}$;
        \item новое дерево $h_m$;
        \item значение общей функции потерь.
    \end{itemize}

    \mlStatement

    Основной наблюдаемый эффект --- последовательное уменьшение ошибки ансамбля за счёт
    добавления моделей, приближающих антиградиент функции потерь.

    \begin{columns}[T,onlytextwidth]
        \column{0.50\textwidth}
        \centering
        \safeimg[width=0.95\linewidth]{img/boosting_steps.pdf}
        \column{0.50\textwidth}
        \centering
        \safeimg[width=0.95\linewidth]{img/boosting_loss.pdf}
    \end{columns}

    % Manim: animations/08_gradient_boosting.mp4
\end{frame}

% ============================================================
% Итоги
% ============================================================

\begin{frame}{Итоги}
    \begin{enumerate}
        \item Обучение параметрической модели формулируется как задача
        \[
            \hat\theta\in\arg\min_{\theta}L_D(\theta).
        \]

        \item Для гладких функций потерь естественно возникают функции нескольких
        переменных, частные производные, градиент, линии уровня и численные методы
        поиска минимума.

        \item Linear Regression и Logistic Regression различаются видом модели и
        функции потерь, но сохраняют общую постановку оптимизации.

        \item Decision Tree использует дискретное жадное построение структуры,
        тогда как Gradient Boosting возвращает идею антиградиента на уровне
        последовательного построения функции.

        \item PCA, Monte Carlo и Ant Colony Optimization иллюстрируют использование
        линейной алгебры, вероятностных методов и метаэвристической оптимизации.
    \end{enumerate}

    \mlConsequence

    Математический анализ возникает в машинном обучении не как отдельный набор формул,
    а как естественный язык постановки и решения задач обучения.
\end{frame}

\begin{frame}{Основные источники}
    \small
    \begin{itemize}
        \item Школа анализа данных Яндекса. \textit{Учебник по машинному обучению}.
        \item Броиловский А. \textit{Обзор задач машинного обучения}. ШАД СПб, 2026.
        \item Овчаров С. \textit{Линейная регрессия и основы оптимизации}. ШАД СПб, 2026.
        \item Овчаров С. \textit{Логистическая регрессия и классификация}. ШАД СПб, 2026.
        \item Zhang A., Lipton Z. C., Li M., Smola A. J.
        \textit{Dive into Deep Learning: Optimization Algorithms}.
        \item Материалы проекта \textit{«Математический анализ в методах машинного обучения»}.
    \end{itemize}

    \mlRemark

    В презентации используются собственные вычислительные эксперименты и визуализации,
    построенные по результатам программной реализации рассматриваемых методов.
\end{frame}

\end{document}
